\documentclass[12pt,a4paper]{article}

\usepackage[utf8]{inputenc}
\usepackage[T1]{fontenc}
\usepackage[english]{babel}
\usepackage{amsmath,amssymb}
\usepackage{siunitx}
\usepackage{geometry}
\usepackage{booktabs}
\usepackage{array}
\usepackage{hyperref}

\geometry{margin=2.5cm}
\sisetup{output-decimal-marker = {.}, per-mode = symbol}
\hypersetup{colorlinks=true,linkcolor=blue,urlcolor=blue}

\title{Optical Sensing Circuit Component Calculations\\
\large Glucometer Project}
\author{}
\date{\today}

\begin{document}
\maketitle

\section{Objective}
This document presents the calculations used to size the LED current-limiting
resistor and the transimpedance amplifier feedback network for two emitter
and photodiode combinations:
\begin{itemize}
  \item TSAL6100 LED with BPV10NF photodiode.
  \item SFH 4544 LED with QSB34GR photodiode.
\end{itemize}

The supplied PDF schematic is used as a visual reference for the circuit
architecture and the components shown. The numerical results below are
calculated from the parameter values specified for each combination.

\section{Design Equations}

\subsection{LED Current-Limiting Resistor}
The resistor connected in series with the LED is calculated as:
\begin{equation}
  R_3 = \frac{V_A - V_L}{I_{LED}},
  \label{eq:rled}
\end{equation}
where $V_A$ is the supply voltage, $V_L$ is the LED forward voltage, and
$I_{LED}$ is the desired LED current.

\subsection{Total Input Capacitance}
The total capacitance considered at the amplifier input is:
\begin{equation}
  C_{in} = C_{pd} + C_{inA},
  \label{eq:cin}
\end{equation}
where
\begin{equation}
  C_{inA} = C_{diff} + C_{cm}.
\end{equation}
Here, $C_{pd}$ is the photodiode junction capacitance, $C_{diff}$ is the
operational amplifier differential input capacitance, and $C_{cm}$ is the
common-mode capacitance at the inverting input. The calculations use the
provided combined value $C_{inA}$.

\subsection{Feedback Resistor}
For a maximum photodiode current $I_{pd}$ and a target output voltage $V_o$,
the feedback resistor is estimated using:
\begin{equation}
  R_4 = \frac{V_o}{I_{pd}}.
  \label{eq:rf}
\end{equation}

\subsection{Compensation Capacitor}
The estimated feedback capacitor for the stated stability condition is:
\begin{equation}
  C_1 = \sqrt{\frac{C_{in}}{2\pi R_4\,GBW}},
  \label{eq:cf}
\end{equation}
where $GBW$ is the operational amplifier gain-bandwidth product.

\subsection{Transfer Function}
Using the transfer-function expression provided:
\begin{equation}
  V_{out}(s) =
  I_{pd}\left(\frac{R_4}{1+sR_4C_1}\right).
  \label{eq:transfer}
\end{equation}
This is the magnitude-form expression specified for the transimpedance
stage. In an inverting transimpedance configuration, the output sign depends
on the current convention and photodiode connection.

\section{Case 1: TSAL6100 and BPV10NF}

\subsection{Design Parameters}
\begin{table}[h!]
\centering
\begin{tabular}{>{\raggedright\arraybackslash}p{7cm}r}
\toprule
Parameter & Value \\
\midrule
Supply voltage, $V_A$ & \SI{3.3}{\volt}\\
LED forward voltage, $V_L$ & \SI{1.35}{\volt}\\
LED current, $I_{LED}$ & \SI{30}{\milli\ampere}\\
Photodiode capacitance, $C_{pd}$ & \SI{11}{\pico\farad}\\
Amplifier input capacitance, $C_{inA}$ & \SI{25}{\pico\farad}\\
Gain-bandwidth product, $GBW$ & \SI{2.5}{\mega\hertz}\\
Target output voltage, $V_o$ & \SI{1.65}{\volt}\\
Maximum photodiode current, $I_{pd}$ & \SI{60}{\micro\ampere}\\
\bottomrule
\end{tabular}
\caption{Parameters used for Case 1.}
\end{table}

\subsection{Calculation of $R_3$}
Substituting the values into Equation~\eqref{eq:rled}:
\begin{align}
  R_3 &= \frac{\SI{3.3}{\volt}-\SI{1.35}{\volt}}
               {\SI{30}{\milli\ampere}}\\
      &= \frac{\SI{1.95}{\volt}}{\SI{0.03}{\ampere}}\\
      &= \boxed{\SI{65}{\ohm}}.
\end{align}
A nearby standard resistor value is \SI{68}{\ohm}. With this value, the
approximate LED current is:
\begin{equation}
  I_{LED} = \frac{\SI{3.3}{\volt}-\SI{1.35}{\volt}}
                 {\SI{68}{\ohm}}
           \approx \SI{28.68}{\milli\ampere}.
\end{equation}

\subsection{Calculation of $C_{in}$}
\begin{align}
  C_{in} &= C_{pd}+C_{inA}\\
         &= \SI{11}{\pico\farad}+\SI{25}{\pico\farad}\\
         &= \boxed{\SI{36}{\pico\farad}}.
\end{align}

\subsection{Calculation of $R_4$}
\begin{align}
  R_4 &= \frac{\SI{1.65}{\volt}}{\SI{60}{\micro\ampere}}\\
      &= \boxed{\SI{27.5}{\kilo\ohm}}.
\end{align}

\subsection{Calculation of $C_1$}
\begin{align}
 C_1 &=
 \sqrt{\frac{\SI{36e-12}{\farad}}
 {2\pi(\SI{27.5e3}{\ohm})(\SI{2.5e6}{\hertz})}}\\
 &= \boxed{\SI{9.13}{\pico\farad}}.
\end{align}
A nearby standard value is \SI{10}{\pico\farad}.

\section{Case 2: SFH 4544 and QSB34GR}

\subsection{Design Parameters}
\begin{table}[h!]
\centering
\begin{tabular}{>{\raggedright\arraybackslash}p{7cm}r}
\toprule
Parameter & Value \\
\midrule
Supply voltage, $V_A$ & \SI{3.3}{\volt}\\
LED forward voltage, $V_L$ & \SI{1.6}{\volt}\\
LED current, $I_{LED}$ & \SI{100}{\milli\ampere}\\
Photodiode capacitance, $C_{pd}$ & \SI{25}{\pico\farad}\\
Amplifier input capacitance, $C_{inA}$ & \SI{25}{\pico\farad}\\
Gain-bandwidth product, $GBW$ & \SI{2.5}{\mega\hertz}\\
Target output voltage, $V_o$ & \SI{1.65}{\volt}\\
Maximum photodiode current, $I_{pd}$ & \SI{37}{\micro\ampere}\\
\bottomrule
\end{tabular}
\caption{Parameters used for Case 2.}
\end{table}

\subsection{Calculation of $R_3$}
\begin{align}
  R_3 &= \frac{\SI{3.3}{\volt}-\SI{1.6}{\volt}}
               {\SI{100}{\milli\ampere}}\\
      &= \frac{\SI{1.7}{\volt}}{\SI{0.1}{\ampere}}\\
      &= \boxed{\SI{17}{\ohm}}.
\end{align}
Therefore, for the specified values, the result is not \SI{68}{\ohm}.
Using a \SI{68}{\ohm} resistor would result in an estimated current of:
\begin{equation}
  I_{LED} = \frac{\SI{1.7}{\volt}}{\SI{68}{\ohm}}
           \approx \SI{25}{\milli\ampere}.
\end{equation}

\subsection{Calculation of $C_{in}$}
\begin{align}
  C_{in} &= C_{pd}+C_{inA}\\
         &= \SI{25}{\pico\farad}+\SI{25}{\pico\farad}\\
         &= \boxed{\SI{50}{\pico\farad}}.
\end{align}

\subsection{Calculation of $R_4$}
\begin{align}
  R_4 &= \frac{\SI{1.65}{\volt}}{\SI{37}{\micro\ampere}}\\
      &\approx \boxed{\SI{44.59}{\kilo\ohm}}.
\end{align}

\subsection{Calculation of $C_1$}
\begin{align}
 C_1 &=
 \sqrt{\frac{\SI{50e-12}{\farad}}
 {2\pi(\SI{44.5946e3}{\ohm})(\SI{2.5e6}{\hertz})}}\\
 &\approx \boxed{\SI{8.45}{\pico\farad}}.
\end{align}
Nearby standard values include \SI{8.2}{\pico\farad} and
\SI{10}{\pico\farad}; the final choice should be validated by checking
the stability and response of the actual circuit.

\section{Summary of Results}
\begin{table}[h!]
\centering
\begin{tabular}{lcc}
\toprule
Parameter & TSAL6100 / BPV10NF & SFH 4544 / QSB34GR\\
\midrule
Calculated $R_3$ & \SI{65}{\ohm} & \SI{17}{\ohm}\\
Suggested standard $R_3$ & \SI{68}{\ohm} & \SI{17}{\ohm}\\
$C_{in}$ & \SI{36}{\pico\farad} & \SI{50}{\pico\farad}\\
Calculated $R_4$ & \SI{27.5}{\kilo\ohm} & \SI{44.59}{\kilo\ohm}\\
Calculated $C_1$ & \SI{9.13}{\pico\farad} & \SI{8.45}{\pico\farad}\\
\bottomrule
\end{tabular}
\caption{Component calculation summary for both configurations.}
\end{table}

\section{Reference Schematic}
The supplied PDF schematic shows an optical sensing stage using an
LMP7704MA operational amplifier, a TSAL6100 emitter LED, a BPV10NF
photodiode, and feedback components identified as $R_4$ and $C_1$.
The schematic includes reference values of \SI{68}{\ohm} for the LED
resistor, \SI{27}{\kilo\ohm} for $R_4$, and \SI{10}{\pico\farad} for
$C_1$. These values are close to the calculated results for Case 1.

The PDF is used as a reference for the drawn implementation. It does not
confirm the electrical characteristics of the second combination,
SFH 4544/QSB34GR.

\section{Implementation Considerations}
\begin{itemize}
  \item Verify the actual LED forward voltage at the selected operating
        current using the component datasheet.
  \item Ensure that the \SI{3.3}{\volt} supply and the LED driver or
        transistor can provide the required current. In particular,
        \SI{100}{\milli\ampere} must not be assumed to be available from
        an ESP32 GPIO pin.
  \item Select the power rating of $R_3$ based on
        $P=I_{LED}^{2}R_3$, including an appropriate thermal margin.
  \item Check that the maximum photodiode current and output voltage do
        not drive the operational amplifier into saturation.
  \item The equation for $C_1$ provides a design estimate. Stability,
        bandwidth, and transient response should be verified using the
        actual circuit, component datasheets, and measurements.
\end{itemize}

\end{document}
